\begin{figure}[htbp]
\centering
\begin{tikzpicture}[scale=0.95]
    % ========== Row 1: Impulse <-> 1 ==========
    \begin{scope}[yshift=5.5cm]
        % Time domain
        \begin{scope}[xshift=-5cm]
            \node[font=\bfseries, UofSCGarnet] at (1.5, 1.8) {Time Domain};
            \draw[thick,->] (-0.5,0) -- (3.5,0) node[right, font=\small] {$t$};
            \draw[thick,->] (0,-0.3) -- (0,1.5) node[above, font=\small] {$f(t)$};

            % Impulse arrow
            \draw[UofSCGarnet, very thick, ->] (0,0) -- (0,1.2);
            \node[UofSCGarnet, font=\small, right] at (0.15, 0.9) {$\delta(t)$};
            \node[font=\small] at (1.8, -0.5) {Unit Impulse};
        \end{scope}

        % Arrow
        \node[font=\Large] at (0, 0.5) {$\xleftrightarrow{\quad\Laplace\quad}$};

        % s-domain
        \begin{scope}[xshift=5cm]
            \node[font=\bfseries, UofSCAtlantic] at (1.5, 1.8) {$s$-Domain};
            \draw[thick,->] (-0.5,0) -- (3.5,0) node[right, font=\small] {$\sigma$};
            \draw[thick,->] (0,-0.3) -- (0,1.5) node[above, font=\small] {$\jj\omega$};

            % Constant function
            \draw[UofSCAtlantic, very thick] (-0.3, 0.8) -- (3.2, 0.8);
            \node[UofSCAtlantic, font=\small, right] at (3.2, 0.8) {$F(s) = 1$};
            \node[font=\small] at (1.8, -0.5) {No poles or zeros};
        \end{scope}
    \end{scope}

    % ========== Row 2: Step <-> 1/s ==========
    \begin{scope}[yshift=1.8cm]
        % Time domain
        \begin{scope}[xshift=-5cm]
            \draw[thick,->] (-0.5,0) -- (3.5,0) node[right, font=\small] {$t$};
            \draw[thick,->] (0,-0.3) -- (0,1.5) node[above, font=\small] {$f(t)$};

            % Step function
            \draw[UofSCGarnet, very thick] (-0.3, 0) -- (0, 0);
            \draw[UofSCGarnet, very thick] (0, 0) -- (0, 1);
            \draw[UofSCGarnet, very thick] (0, 1) -- (3.2, 1);
            \node[font=\small] at (1.8, 1.3) {$u(t)$};
            \node[font=\small] at (1.8, -0.5) {Unit Step};
        \end{scope}

        % Arrow
        \node[font=\Large] at (0, 0.5) {$\xleftrightarrow{\quad\Laplace\quad}$};

        % s-domain
        \begin{scope}[xshift=5cm]
            \draw[thick,->] (-0.5,0) -- (3.5,0) node[right, font=\small] {$\sigma$};
            \draw[thick,->] (0,-0.3) -- (0,1.5) node[above, font=\small] {$\jj\omega$};

            % Pole at origin
            \node[UofSCGarnet, font=\Large] at (0, 0) {$\times$};
            \node[UofSCAtlantic, font=\small, right] at (2, 0.8) {$F(s) = \dfrac{1}{s}$};
            \node[font=\small] at (1.8, -0.5) {Pole at $s = 0$};
        \end{scope}
    \end{scope}

    % ========== Row 3: Exponential <-> 1/(s+a) ==========
    \begin{scope}[yshift=-1.9cm]
        % Time domain
        \begin{scope}[xshift=-5cm]
            \draw[thick,->] (-0.5,0) -- (3.5,0) node[right, font=\small] {$t$};
            \draw[thick,->] (0,-0.3) -- (0,1.5) node[above, font=\small] {$f(t)$};

            % Exponential decay
            \draw[UofSCHorseshoe, very thick, domain=0:3.2, samples=50] plot (\x, {1.2*exp(-1.2*\x)});
            \node[font=\small] at (1.8, 1.3) {$e^{-at}u(t)$};
            \node[font=\small] at (1.8, -0.5) {Decaying Exponential};

            % Time constant marker
            \draw[UofSCCongaree, dashed, thin] (0.833, 0) -- (0.833, 0.44);
            \node[UofSCCongaree, font=\tiny, below] at (0.833, 0) {$\tau = 1/a$};
        \end{scope}

        % Arrow
        \node[font=\Large] at (0, 0.5) {$\xleftrightarrow{\quad\Laplace\quad}$};

        % s-domain
        \begin{scope}[xshift=5cm]
            \draw[thick,->] (-1,0) -- (3.5,0) node[right, font=\small] {$\sigma$};
            \draw[thick,->] (0,-0.3) -- (0,1.5) node[above, font=\small] {$\jj\omega$};

            % Pole at s = -a
            \node[UofSCGarnet, font=\Large] at (-0.7, 0) {$\times$};
            \node[UofSCGarnet, font=\tiny, below] at (-0.7, -0.15) {$-a$};
            \node[UofSCAtlantic, font=\small, right] at (1.5, 0.8) {$F(s) = \dfrac{1}{s+a}$};
            \node[font=\small] at (1.8, -0.5) {Pole at $s = -a$};

            % ROC indication
            \fill[UofSCHorseshoe!15] (-0.7, -0.25) rectangle (3.3, 1.4);
            \draw[UofSCCongaree, dashed] (-0.7, -0.25) -- (-0.7, 1.4);
        \end{scope}
    \end{scope}

    % ========== Row 4: Sine <-> omega/(s^2+omega^2) ==========
    \begin{scope}[yshift=-5.6cm]
        % Time domain
        \begin{scope}[xshift=-5cm]
            \draw[thick,->] (-0.5,0) -- (3.5,0) node[right, font=\small] {$t$};
            \draw[thick,->] (0,-1.2) -- (0,1.5) node[above, font=\small] {$f(t)$};

            % Sine wave
            \draw[UofSCRose, very thick, domain=0:3.2, samples=100] plot (\x, {sin(deg(3*\x))});
            \node[font=\small] at (1.8, 1.5) {$\sin(\omega t) \cdot u(t)$};
            \node[font=\small] at (1.8, -1.5) {Sinusoidal};
        \end{scope}

        % Arrow
        \node[font=\Large] at (0, 0) {$\xleftrightarrow{\quad\Laplace\quad}$};

        % s-domain
        \begin{scope}[xshift=5cm]
            \draw[thick,->] (-0.5,0) -- (3.5,0) node[right, font=\small] {$\sigma$};
            \draw[thick,->] (0,-1.2) -- (0,1.5) node[above, font=\small] {$\jj\omega$};

            % Poles on imaginary axis
            \node[UofSCGarnet, font=\Large] at (0, 0.8) {$\times$};
            \node[UofSCGarnet, font=\Large] at (0, -0.8) {$\times$};
            \node[UofSCGarnet, font=\tiny, right] at (0.15, 0.8) {$\jj\omega$};
            \node[UofSCGarnet, font=\tiny, right] at (0.15, -0.8) {$-\jj\omega$};

            \node[UofSCAtlantic, font=\small, right] at (1.2, 0.3) {$F(s) = \dfrac{\omega}{s^2 + \omega^2}$};
            \node[font=\small] at (1.8, -1.5) {Poles at $s = \pm\jj\omega$};
        \end{scope}
    \end{scope}

    % Dividing lines
    \draw[gray!50, thick] (-8, 4) -- (8, 4);
    \draw[gray!50, thick] (-8, 0.3) -- (8, 0.3);
    \draw[gray!50, thick] (-8, -3.4) -- (8, -3.4);
\end{tikzpicture}

\caption{Visual comparison of fundamental Laplace transform pairs showing the correspondence between time-domain signals (left, Garnet/Horseshoe/Rose) and their $s$-domain representations (right, Atlantic). From top to bottom: (1) the unit impulse $\delta(t)$ transforms to constant 1 with no poles; (2) the unit step $u(t)$ transforms to $1/s$ with a pole at the origin; (3) the decaying exponential $e^{-at}u(t)$ transforms to $1/(s+a)$ with a pole at $s=-a$ (the shaded region shows the ROC); (4) the sine wave $\sin(\omega t)u(t)$ transforms to $\omega/(s^2+\omega^2)$ with complex conjugate poles on the imaginary axis. Pole locations directly determine the time-domain behavior.}
\label{fig:laplace_time_domain_pairs}
\end{figure}
